

\documentclass{article}
\usepackage{amsmath,amssymb,amsthm,enumitem}

\newtheorem{theorem}{Theorem}
\newtheorem{lemma}[theorem]{Lemma}
\newtheorem{proposition}[theorem]{Proposition}
\newtheorem{corollary}[theorem]{Corollary}
\theoremstyle{remark}
\newtheorem{remark}[theorem]{Remark}

\begin{document}

\title{Cohomological descent for the inclusion of \(B_P^{\rm op}\) into \(A_P^{\rm op}\)}
\author{}
\maketitle

\section{Answer}

Yes.  For every monoid \(P\) and every presheaf of abelian groups
\(F\) on \(B_P^{\rm op}\), the canonical morphism
\[
F \longrightarrow \mathrm R\pi_\ast\pi^{-1}F
\]
is a quasi-isomorphism.  In fact, under the hypotheses stated in the
question, the direct image functor \(\pi_\ast\) is exact and the
underived adjunction unit
\[
F \longrightarrow \pi_\ast\pi^{-1}F
\]
is an isomorphism.  Hence the derived unit is an isomorphism as well.

The proof is essentially topos-theoretic.  The monoidal definitions
enter only through the asserted property that the inclusion
\(B_P^{\rm op}\subset A_P^{\rm op}\) is left exact as a functor of
sites.  Since \(B_P^{\rm op}\) is given the trivial topology, this
implies that every covering sieve of an object of \(B_P^{\rm op}\) in
\(A_P^{\rm op}\) already contains the identity.  This makes the
restriction functor exact and forces sheafification not to change the
values of presheaves on objects of \(B_P^{\rm op}\).

\section{Notation}

Let
\[
A=A_P,\qquad B=B_P,
\]
and let
\[
j:B\hookrightarrow A
\]
be the fully faithful inclusion.  The sites in the problem are
\[
\mathcal C=A^{\rm op},\qquad \mathcal D=B^{\rm op},
\]
with \(\mathcal C\) equipped with the birhizomic topology and
\(\mathcal D\) equipped with the trivial topology.  We write
\[
i:\mathcal D\hookrightarrow \mathcal C
\]
for the induced inclusion \(i=j^{\rm op}\).

A presheaf on \(\mathcal D=B^{\rm op}\) is the same thing as a covariant
functor
\[
F:B\longrightarrow \mathbf{Ab}.
\]
The left Kan extension used in the question is the left Kan extension
along \(j\):
\[
\operatorname{Lan}_j F:A\longrightarrow \mathbf{Ab},
\]
regarded as a presheaf on \(A^{\rm op}\).  Sheafification on
\(\mathcal C=A^{\rm op}\) is denoted by \(a\).  Thus
\[
\pi^{-1}F = a(\operatorname{Lan}_j F).
\]

The right adjoint \(\pi_\ast\) is restriction along \(j\):
\[
\pi_\ast G = G\circ j
\]
for a sheaf \(G\) on \(\mathcal C\).  Indeed, for \(F\in
\operatorname{PSh}(\mathcal D)\) and \(G\in \operatorname{Sh}(\mathcal C)\),
\[
\operatorname{Hom}_{\operatorname{Sh}(\mathcal C)}(a\operatorname{Lan}_jF,G)
\simeq
\operatorname{Hom}_{\operatorname{PSh}(\mathcal C)}(\operatorname{Lan}_jF,G)
\simeq
\operatorname{Hom}_{\operatorname{PSh}(\mathcal D)}(F,G\circ j).
\]

\section{A general topos-theoretic lemma}

We first prove a general statement which directly implies the desired
conclusion.

\begin{lemma}\label{lemma:discrete}
Let \(i:\mathcal D\hookrightarrow \mathcal C\) be a fully faithful
inclusion of sites.  Assume that \(\mathcal D\) has the trivial
topology, and that \(i\) is continuous.  Then for every object
\(U\in\mathcal D\) and every covering sieve \(R\) on \(i(U)\) in
\(\mathcal C\), the identity morphism \(\mathrm{id}_{i(U)}\) belongs to
\(R\).
\end{lemma}

\begin{proof}
Continuity of \(i\) means that for every covering sieve \(R\) on
\(i(U)\), the pulled-back sieve
\[
i^{-1}R=\{\,f:V\to U\text{ in }\mathcal D \mid i(f)\in R\,\}
\]
is a covering sieve in \(\mathcal D\).  Since the topology on
\(\mathcal D\) is trivial, the only covering sieve on \(U\) is the
maximal sieve.  Hence \(i^{-1}R\) is maximal.  In particular,
\(\mathrm{id}_U\in i^{-1}R\), and therefore
\[
i(\mathrm{id}_U)=\mathrm{id}_{i(U)}\in R.
\]
\end{proof}

\begin{remark}
In the terminology used in the problem, the assertion that the
inclusion \(B_P^{\rm op}\subset A_P^{\rm op}\) is left exact is made
as a functor of sites; in particular it includes continuity.  Thus
Lemma~\ref{lemma:discrete} applies to
\(\mathcal D=B_P^{\rm op}\subset \mathcal C=A_P^{\rm op}\).
\end{remark}

\begin{lemma}\label{lemma:sheafification-value}
Let \(\mathcal C\) be a site and let \(c\in\mathcal C\) be an object
such that every covering sieve on \(c\) contains
\(\mathrm{id}_c\).  Then for every presheaf of abelian groups
\(G\) on \(\mathcal C\), the canonical map
\[
G(c)\longrightarrow (aG)(c)
\]
to the associated sheaf is an isomorphism.
\end{lemma}

\begin{proof}
Use the usual plus construction for sheafification of abelian
presheaves.  For a presheaf \(G\), define
\[
G^+(c)=\varinjlim_{R} \operatorname{Hom}(R,G),
\]
where \(R\) runs over covering sieves on \(c\), and
\(\operatorname{Hom}(R,G)\) is the group of matching families for
\(G\) over \(R\).

If \(R\) is a covering sieve on \(c\), then by hypothesis
\(\mathrm{id}_c\in R\).  The category of elements of \(R\) then has a
terminal object, namely \(\mathrm{id}_c:c\to c\).  Therefore a matching
family over \(R\) is uniquely determined by its component along
\(\mathrm{id}_c\), and every element of \(G(c)\) defines a matching
family.  Hence the restriction map
\[
G(c)\xrightarrow{\sim}\operatorname{Hom}(R,G)
\]
is an isomorphism.  These isomorphisms are compatible with the
transition maps induced by inclusions of sieves, so
\[
G^+(c)\cong G(c).
\]
Sheafification is obtained by applying the plus construction twice, so
\((aG)(c)\cong G(c)\).  Under this identification the canonical map is
the identity.
\end{proof}

\begin{lemma}\label{lemma:pi-star-exact}
Let \(i:\mathcal D\hookrightarrow \mathcal C\) be a fully faithful
inclusion of sites satisfying the conclusion of
Lemma~\ref{lemma:discrete}: every covering sieve on an object of
\(\mathcal D\) contains the identity.  Let
\[
\pi^{-1}:\operatorname{Sh}(\mathcal D)\to \operatorname{Sh}(\mathcal C)
\]
be the functor given by left Kan extension followed by sheafification,
and let \(\pi_\ast\) be its right adjoint.  Then \(\pi_\ast\) is exact.
\end{lemma}

\begin{proof}
As observed above, \(\pi_\ast\) is restriction along the inclusion
\(\mathcal D\subset\mathcal C\).  We must show that restriction from
sheaves on \(\mathcal C\) to presheaves on \(\mathcal D\) preserves
short exact sequences.

Finite limits of sheaves of abelian groups are computed objectwise, so
restriction preserves kernels.  It remains to prove that restriction
preserves epimorphisms.

Let \(\varphi:G\to H\) be an epimorphism of sheaves on \(\mathcal C\),
and let \(U\in\mathcal D\).  Take a section \(s\in H(U)\).  By the
standard characterization of epimorphisms of sheaves, the sieve
\[
R_s=\{\,f:V\to U \mid s|_V \text{ lifts to } G(V)\,\}
\]
is a covering sieve on \(U\); see for instance the Stacks Project,
Lemma 00ZG.  By hypothesis, \(\mathrm{id}_U\in R_s\).  Hence \(s\)
itself lifts to a section of \(G(U)\).  Thus
\[
G(U)\longrightarrow H(U)
\]
is surjective for every \(U\in\mathcal D\), and restriction preserves
epis.  Therefore restriction is exact.
\end{proof}

\begin{lemma}\label{lemma:underived-unit}
Let \(i:\mathcal D\hookrightarrow \mathcal C\) be as in
Lemma~\ref{lemma:pi-star-exact}.  For every presheaf of abelian groups
\(F\) on \(\mathcal D\), the underived adjunction unit
\[
F\longrightarrow \pi_\ast\pi^{-1}F
\]
is an isomorphism.
\end{lemma}

\begin{proof}
Let \(j:\mathcal D^{\rm op}\hookrightarrow \mathcal C^{\rm op}\) be the
opposite inclusion.  The presheaf \(F\) on \(\mathcal D\) is a functor
\(\mathcal D^{\rm op}\to\mathbf{Ab}\), and the relevant left Kan
extension is \(\operatorname{Lan}_j F\).

Since \(j\) is fully faithful, left Kan extension along \(j\) restricts
back to \(F\).  Concretely, for \(U\in\mathcal D\),
\[
(\operatorname{Lan}_jF)(U)
=
\varinjlim_{(V\in\mathcal D,\ V\to U)} F(V).
\]
The indexing category has a terminal object, namely
\((U,\mathrm{id}_U)\), so the colimit is \(F(U)\).  Thus
\[
(\operatorname{Lan}_jF)|_{\mathcal D}\cong F.
\]

Now sheafify.  By Lemma~\ref{lemma:sheafification-value}, sheafification
does not change the value at any object of \(\mathcal D\), because every
covering sieve on such an object contains the identity.  Hence
\[
(\pi^{-1}F)|_{\mathcal D}
=
(a\operatorname{Lan}_jF)|_{\mathcal D}
\cong
(\operatorname{Lan}_jF)|_{\mathcal D}
\cong
F.
\]
Since \(\pi_\ast\) is restriction, this gives
\[
\pi_\ast\pi^{-1}F\cong F.
\]
The adjunction unit is precisely this composite isomorphism.
\end{proof}

\begin{proposition}\label{prop:general}
Let \(i:\mathcal D\hookrightarrow \mathcal C\) be a fully faithful
inclusion of sites satisfying the hypotheses of
Lemma~\ref{lemma:pi-star-exact}.  Then for every presheaf of abelian
groups \(F\) on \(\mathcal D\), the morphism
\[
F\longrightarrow \mathrm R\pi_\ast\pi^{-1}F
\]
is a quasi-isomorphism.
\end{proposition}

\begin{proof}
By Lemma~\ref{lemma:pi-star-exact}, \(\pi_\ast\) is exact.  Therefore
its right derived functor is just \(\pi_\ast\) itself:
\[
\mathrm R\pi_\ast=\pi_\ast.
\]
Thus
\[
\mathrm R\pi_\ast\pi^{-1}F
=
\pi_\ast\pi^{-1}F.
\]
By Lemma~\ref{lemma:underived-unit}, the underived unit
\(F\to\pi_\ast\pi^{-1}F\) is an isomorphism.  Hence the derived unit is
a quasi-isomorphism.
\end{proof}

\section{Application to the birhizomic site}

We now apply Proposition~\ref{prop:general} to the situation of the
question.

Recall that
\[
\mathcal C=A_P^{\rm op},\qquad \mathcal D=B_P^{\rm op}.
\]
The topology on \(\mathcal D\) is trivial, and the inclusion
\[
i:\mathcal D\hookrightarrow \mathcal C
\]
is fully faithful.  The problem states that this inclusion is left
exact.  As a functor of sites, this means in particular that it is
continuous.  Therefore Lemma~\ref{lemma:discrete} applies: every
birhizomic covering sieve on an object of \(B_P^{\rm op}\) contains the
identity.

Consequently the hypotheses of Proposition~\ref{prop:general} are
satisfied.  Hence for every presheaf of abelian groups \(F\) on
\(B_P^{\rm op}\), the canonical morphism
\[
F\longrightarrow \mathrm R\pi_\ast\pi^{-1}F
\]
is a quasi-isomorphism.

\begin{theorem}
For every monoid \(P\) and every presheaf of abelian groups \(F\) on
\(B_P^{\rm op}\), the morphism
\[
F \longrightarrow \mathrm R\pi_\ast\pi^{-1}F
\]
is a quasi-isomorphism.
\end{theorem}

\begin{proof}
This is exactly Proposition~\ref{prop:general} applied to
\(\mathcal C=A_P^{\rm op}\) and \(\mathcal D=B_P^{\rm op}\), using the
continuity consequence of the stated left exactness of the inclusion.
\end{proof}

\begin{remark}
The proof shows more: the functor \(\pi_\ast\) is exact and the
underived unit \(F\to\pi_\ast\pi^{-1}F\) is an isomorphism.  Thus the
derived statement is not merely a vanishing theorem; it follows from an
actual equality at the abelian level.
\end{remark}

\end{document}